\documentclass[11pt,a4paper]{article}
\usepackage[utf8]{inputenc}
\usepackage[T1]{fontenc}
\usepackage[spanish,es-minimal,es-noshorthands]{babel}
\usepackage{lmodern}
\linespread{1.22}
\usepackage[a4paper,top=2.4cm,bottom=2.2cm,left=1.5cm,right=1.5cm,headheight=1.1cm,headsep=0.5cm,footskip=1cm]{geometry}
\usepackage{amsmath,amssymb,cancel}
\usepackage{icomma}
\usepackage{graphicx}
\usepackage[table]{xcolor}
\usepackage{tikz}\usetikzlibrary{calc}
\usepackage{pgfplots}\pgfplotsset{compat=1.17}
\usepackage{fontawesome5}
\usepackage{ccicons}
\usepackage{fancyhdr}
\usepackage{enumitem}
\usepackage{booktabs}
\usepackage{tabularx}
\usepackage{needspace}
\usepackage[most]{tcolorbox}
\usepackage{titlesec}
\usepackage[hidelinks]{hyperref}
\hypersetup{pdftitle={Apuntes de teoría. Unidad 2: Álgebra. 1.º de Bachillerato CCSS},pdfsubject={Matemáticas Aplicadas a las Ciencias Sociales I},pdfkeywords={polinomios, identidades notables, regla de Ruffini, teorema del resto, factorización, fracciones algebraicas, ecuaciones, inecuaciones}}

% colors de la web (FentMates, bloque de Álgebra: azul; es manté el nom fmgreen dels apunts de les altres unitats)
\definecolor{fmnavy}{HTML}{1B2748}
\definecolor{fmgreen}{HTML}{3D6AA0}
\definecolor{fmgreenl}{HTML}{4E81BD}
\definecolor{fmyellow}{HTML}{FFC94D}
\definecolor{fmgold}{HTML}{E9A919}
\definecolor{fmaccent}{HTML}{FF7A50}
\definecolor{fmteal}{HTML}{2BB3A3}
\definecolor{fmdef}{HTML}{1D4F73}\definecolor{fmdefbg}{HTML}{EEF4F9}
\definecolor{fmprop}{HTML}{3B4252}\definecolor{fmpropbg}{HTML}{F3F4F7}
\definecolor{fmex}{HTML}{A3203A}\definecolor{fmexbg}{HTML}{FBF0F2}
\definecolor{fmmuted}{HTML}{5D6685}\definecolor{fmball0}{HTML}{E2483D}\definecolor{fmball1}{HTML}{2F6FD6}
\definecolor{fmsurface}{HTML}{EEF1F8}
\definecolor{fmA}{HTML}{3B82F6}\definecolor{fmB}{HTML}{F97316}\definecolor{fmR}{HTML}{E11D48}\definecolor{fmK}{HTML}{16A34A}\definecolor{fmV}{HTML}{7C3AED}

% blocs que no es poden partir: es mesuren i es reserva exactament el seu espai (amb el títol o la capçalera del requadre)
\newsavebox{\fmkb}
\newcommand{\fmmeasure}[2]{\setbox\fmkb=\vbox{\hsize=\dimexpr\linewidth-#1\relax\linewidth=\hsize #2}}
\newcommand{\fmblock}[1]{\fmmeasure{0pt}{#1}\Needspace{\dimexpr\ht\fmkb+\dp\fmkb\relax}#1}
\newcommand{\fmheadblock}[2]{\fmmeasure{0pt}{#2}\Needspace{\dimexpr\ht\fmkb+\dp\fmkb+2.2cm\relax}#1#2}
\newcommand{\fmboxneed}[1]{\fmmeasure{9mm}{#1}\Needspace{\dimexpr\ht\fmkb+\dp\fmkb+1.6cm\relax}}
\newif\iffmstack\newif\iffmfit
\newcommand{\fmfits}[1]{\par\ifdim\dimexpr\pagegoal-\pagetotal\relax<\dimexpr\ht\fmkb+\dp\fmkb+#1\relax\global\fmfitfalse\else\global\fmfittrue\fi}
\newcommand{\fmtwo}[2]{\fmmeasure{0pt}{#1}\fmfits{0pt}\iffmfit #1\else #2\fi}
\newcommand{\fmheadtwo}[3]{\fmmeasure{0pt}{#2}\fmfits{2.2cm}\iffmfit #1#2\else\Needspace{\dimexpr1.6cm+7\baselineskip\relax}#1#3\fi}
\newcommand{\fmtwoin}[2]{\iffmstack #2\else #1\fi}
\newcommand{\fmboxneedalt}[1]{\fmmeasure{9mm}{#1}\fmfits{1.6cm}\iffmfit\global\fmstackfalse\else\global\fmstacktrue\Needspace{8\baselineskip}\fi}
\newcommand{\fmheadboxalt}[2]{\fmmeasure{9mm}{#2}\fmfits{3.4cm}\iffmfit\global\fmstackfalse\else\global\fmstacktrue\Needspace{\dimexpr2.4cm+5\baselineskip\relax}\fi#1}
\newcommand{\fmheadbox}[2]{\fmmeasure{9mm}{#2}\Needspace{\dimexpr\ht\fmkb+\dp\fmkb+3.4cm\relax}#1}
\newcommand{\lt}{<}\newcommand{\gt}{>}
% fórmula en línia que no es parteix mai; si no cap en una línia, es posa centrada
% fórmula en línia; si ocupa més del 80 % de la línia, va centrada (amb la puntuació que la segueix)
\makeatletter
\newcommand{\fmm}[1]{\setbox0=\hbox{$#1$}\ifdim\wd0>0.8\linewidth\def\fmmnext{\fmmdisp{#1}}\else\def\fmmnext{\mbox{\box0}}\fi\fmmnext}
\def\fmmdisp#1{\@ifnextchar.{\fmmdispp{#1}}{\@ifnextchar,{\fmmdispc{#1}}{\@ifnextchar;{\fmmdisps{#1}}{\fmdisp{#1}}}}}
\def\fmmdispp#1.{\fmdisp{#1\,.}}\def\fmmdispc#1,{\fmdisp{#1\,,}}\def\fmmdisps#1;{\fmdisp{#1\,;}}
\makeatother
% fórmula centrada; si és més ampla que la línia, s'encongeix fins que hi cap
\newcommand{\fmdisp}[1]{\[\setbox2=\hbox{$\displaystyle #1$}\ifdim\wd2>\linewidth\resizebox{\linewidth}{!}{\box2}\else\box2\fi\]}
\DeclareUnicodeCharacter{2212}{-}
\DeclareUnicodeCharacter{2225}{\ensuremath{\parallel}}\DeclareUnicodeCharacter{22A5}{\ensuremath{\perp}}\DeclareUnicodeCharacter{2226}{\ensuremath{\nparallel}}
\DeclareUnicodeCharacter{207A}{\ensuremath{^{+}}}\DeclareUnicodeCharacter{207B}{\ensuremath{^{-}}}
\setlength{\parindent}{0pt}
\setlength{\parskip}{12pt}
\setlength{\lineskip}{5pt}\setlength{\lineskiplimit}{5pt}  % més aire al voltant de les fraccions en línia
\setlength{\emergencystretch}{3em}
\setlist{itemsep=4pt,topsep=4pt}
\renewcommand{\thesection}{2.\arabic{section}}
\titleformat{\section}{\color{fmnavy}\Large\bfseries}{\colorbox{fmdef}{\color{white}\strut\ \thesection\ }}{0.6em}{}[{\color{fmdef}\titlerule[1.2pt]}]
\titlespacing*{\section}{0pt}{10pt}{4pt}
\titleformat*{\subsection}{\color{fmgreen}\large\bfseries}
\titlespacing*{\subsection}{0pt}{6pt plus 2pt}{0pt}
\allowdisplaybreaks

% dibuix de la targeta de la unitat a la web (illustrations.js), en blanc
\newcommand{\fmcard}{\begin{tikzpicture}[x=0.062cm,y=-0.062cm,white,line width=3pt,line cap=round]
  \begin{scope}[transparency group,opacity=0.15]
  \node[anchor=base west,inner sep=0pt,font=\fontsize{26}{26}\selectfont] at (14,44) {$P(x)=x^2+3x-4$};
  \draw[line width=2.4pt] (40,62) rectangle (84,106) (88,62) rectangle (106,106);
  \fill[opacity=0.3] (40,62) rectangle (84,106);
  \node[anchor=base west,inner sep=0pt,font=\fontsize{17}{17}\selectfont] at (52,90) {$x^2$};
  \node[anchor=base west,inner sep=0pt,font=\fontsize{17}{17}\selectfont] at (91,90) {$x$};
  \draw[line width=3.2pt] (140,106) .. controls (160,106) and (166,66) .. (186,66) .. controls (206,66) and (214,100) .. (226,74);
  \draw[line width=2pt,opacity=0.6] (134,106)--(230,106);
  \end{scope}
\end{tikzpicture}}
\newcommand{\fmbrand}{\raisebox{-0.25\height}{\includegraphics[height=0.5cm]{logo-fentmates.pdf}}}

% capçalera i peu
\pagestyle{fancy}\fancyhf{}
\fancyhead[L]{\small\fmbrand}
\fancyhead[R]{\small\color{fmmuted}1.º Bachillerato CCSS · Unidad 2: Álgebra}
\fancyfoot[L]{\scriptsize\color{fmmuted}\href{https://creativecommons.org/licenses/by-nc-sa/4.0/deed.ca}{\ccbyncsa}\ \ FentMates · fentmates.org · CC BY-NC-SA 4.0}
\fancyfoot[R]{\small\color{fmmuted}\thepage}
\renewcommand{\headrulewidth}{0.4pt}\renewcommand{\headrule}{\color{fmmuted!40}\hrule height \headrulewidth}
\fancypagestyle{primera}{\fancyhf{}\renewcommand{\headrulewidth}{0pt}
  \fancyfoot[L]{\scriptsize\color{fmmuted}\href{https://creativecommons.org/licenses/by-nc-sa/4.0/deed.ca}{\ccbyncsa}\ \ FentMates · fentmates.org · CC BY-NC-SA 4.0}
  \fancyfoot[R]{\small\color{fmmuted}\thepage}}

% caixes (mateixos colors que la web)
% quan una caixa es parteix entre dues pàgines, la vora del tall és discontínua
\tcbset{fmbase/.style={enhanced,breakable,arc=3mm,boxrule=1pt,left=3.5mm,right=3.5mm,top=2mm,bottom=2.5mm,before skip=10pt,after skip=10pt,fonttitle=\bfseries,
  frame style={fill=tcbcolframe},
  overlay first={\fill[tcbcolback] ([xshift=1pt]frame.south west) rectangle ([xshift=-1pt,yshift=1.2pt]frame.south east);
    \draw[tcbcolframe,line width=1pt,dash pattern=on 5pt off 3pt] (frame.south west)--(frame.south east);},
  overlay middle={\fill[tcbcolback] ([xshift=1pt]frame.south west) rectangle ([xshift=-1pt,yshift=1.2pt]frame.south east);
    \fill[tcbcolback] ([xshift=1pt]frame.north west) rectangle ([xshift=-1pt,yshift=-1.2pt]frame.north east);
    \draw[tcbcolframe,line width=1pt,dash pattern=on 5pt off 3pt] (frame.south west)--(frame.south east) (frame.north west)--(frame.north east);},
  overlay last={\fill[tcbcolback] ([xshift=1pt]frame.north west) rectangle ([xshift=-1pt,yshift=-1.2pt]frame.north east);
    \draw[tcbcolframe,line width=1pt,dash pattern=on 5pt off 3pt] (frame.north west)--(frame.north east);}}}
\newtcolorbox{fmdef}[1]{fmbase,colback=fmdefbg,colframe=fmdef,coltitle=white,title={#1}}
\newtcolorbox{fmprop}[1]{fmbase,colback=fmpropbg,colframe=fmprop,coltitle=white,title={#1},halign title=center}
\newtcolorbox{fmex}[1]{fmbase,before upper=\raggedright,colback=fmexbg,colframe=fmex,colbacktitle=fmexbg,coltitle=fmex,title={\faPencil*\ #1},
  titlerule=0.4pt,titlerule style={fmex!40,dashed}}
\newtcolorbox{fmwarn}[1]{fmbase,colback=fmsf!6,colframe=fmsf,coltitle=white,title={#1}}
\newtcolorbox{fmnote}[1][]{enhanced,breakable,colback=white,colframe=white,borderline west={3pt}{0pt}{fmaccent},arc=0pt,left=3mm,right=2mm,top=1mm,bottom=1mm,before skip=6pt,after skip=6pt,before upper=\raggedright,#1}
\newcommand{\fmkey}[1]{\colorbox{fmyellow!45}{\strut #1}}
\newcommand{\fmuline}[1]{\textbf{\color{fmaccent}#1}}
\definecolor{fmok}{HTML}{15803D}\definecolor{fmev}{HTML}{2563EB}\definecolor{fmsf}{HTML}{C2410C}\definecolor{fmsi}{HTML}{7C3AED}\definecolor{fmas}{HTML}{A3203A}
\newcommand{\fmbadge}[2]{\tcbox[on line,colback=#1,colframe=#1,arc=2mm,boxsep=0pt,left=4pt,right=4pt,top=1.5pt,bottom=1.5pt]{\color{white}\footnotesize\bfseries #2}}
\newcommand{\fmchip}[1]{\tcbox[on line,colback=fmsf!10,colframe=fmsf!50,arc=2mm,boxsep=1pt,left=4pt,right=4pt]{#1}}
\newcommand{\fmlevel}[2]{\tcbox[on line,colback=white,colframe=fmmuted!40,arc=2mm,boxsep=2pt,halign=center]{\begin{tabular}{c}{\scriptsize\color{fmmuted}\MakeUppercase{#1}}\\#2\end{tabular}}}
% gràfiques
\pgfplotsset{fmaxis/.style={axis lines=middle,grid=major,grid style={fmmuted!15},axis line style={fmmuted,-latex},
  tick label style={font=\scriptsize,color=fmmuted},clip=true,unbounded coords=jump,restrict y to domain=-200:200,
  every axis plot/.append style={line join=round}},
  fmcurve/.style={fmaccent,line width=1.6pt,smooth}}
\tikzset{fmasy/.style={fmmuted,dashed,line width=0.9pt},fmguide/.style={fmaccent,dotted,line width=0.8pt},
  fmdot/.style={circle,fill=fmaccent,inner sep=0pt,minimum size=5pt},fmopen/.style={circle,draw=fmaccent,fill=white,line width=1pt,inner sep=0pt,minimum size=5pt}}
\newcommand{\fmpau}[1]{\tcbox[on line,colback=fmnavy,colframe=fmnavy,arc=2mm,boxsep=0pt,left=4pt,right=4pt,top=1.5pt,bottom=1.5pt]{\color{white}\scriptsize\bfseries\MakeUppercase{#1}}\ }

\begin{document}
\thispagestyle{primera}
\begin{tikzpicture}[remember picture,overlay]
  \fill[fmgreen] (current page.north west) rectangle ([yshift=-10.1cm]current page.north east);
  \begin{scope}
    \clip (current page.north west) rectangle ([yshift=-10.1cm]current page.north east);
    \node[anchor=north west,inner sep=0pt] at ([xshift=6.6cm,yshift=-1.3cm]current page.north west) {\fmcard};
  \end{scope}
  \node[anchor=north west,inner sep=0pt] at ([xshift=1.5cm,yshift=-1.1cm]current page.north west) {\tcbox[colback=white,colframe=white,arc=2.5mm,boxsep=0pt,left=7pt,right=7pt,top=5pt,bottom=5pt]{\includegraphics[height=0.75cm]{logo-fentmates.pdf}}};
\end{tikzpicture}
\vspace*{0.55cm}
\tcbox[colback=white!15!fmgreen,colframe=white!15!fmgreen,arc=2.5mm,boxsep=0pt,left=6pt,right=6pt,top=3pt,bottom=3pt]{\color{white}\small\bfseries Apuntes de teoría · 1.º Bachillerato · Matemáticas Aplicadas a las Ciencias Sociales I}

\vspace{4pt}
{\raggedright\color{white}\fontsize{25}{30}\selectfont\bfseries{\color{fmyellow}Unidad 2} · Álgebra\par}
\vspace{6pt}
\begin{minipage}{\linewidth}\color{white!88!fmgreen}\small Teoría y ejemplos resueltos paso a paso: polinomios y sus operaciones, la regla de Ruffini y el teorema del resto, la factorización, las fracciones algebraicas (simplificación y operaciones) y las ecuaciones: de primer grado (con paréntesis y denominadores), de segundo grado, bicuadradas y de grado superior, con radicales, racionales, exponenciales y logarítmicas.\end{minipage}\par
% l'índex comença sempre a la mateixa altura, per davall de la franja (com en les unitats 5-10)
\vspace*{\dimexpr 8.6cm-\pagetotal\relax}

\begin{tcolorbox}[colback=fmsurface,colframe=fmsurface,arc=3mm,left=4mm,right=4mm,top=2mm,bottom=2mm]
\small\textbf{Índice}\par\smallskip
\begin{minipage}[t]{0.49\linewidth}\small \mbox{\hyperlink{sec1}{\textbf{\color{fmgreen}2.1}\ \ Polinomios. Operaciones}}\\[3pt]
\mbox{\hyperlink{sec2}{\textbf{\color{fmgreen}2.2}\ \ Ruffini y teorema del resto}}\\[3pt]
\mbox{\hyperlink{sec3}{\textbf{\color{fmgreen}2.3}\ \ Factorización de polinomios}}\\[3pt]
\mbox{\hyperlink{sec4}{\textbf{\color{fmgreen}2.4}\ \ Fracciones algebraicas}}\\[3pt]
\mbox{\hyperlink{sec5}{\textbf{\color{fmgreen}2.5}\ \ Ecuaciones polinómicas}}\\[3pt]\end{minipage}\hfill\begin{minipage}[t]{0.49\linewidth}\small \mbox{\hyperlink{sec6}{\textbf{\color{fmgreen}2.6}\ \ Ecuaciones con radicales}}\\[3pt]
\mbox{\hyperlink{sec7}{\textbf{\color{fmgreen}2.7}\ \ Ecuaciones racionales}}\\[3pt]
\mbox{\hyperlink{sec8}{\textbf{\color{fmgreen}2.8}\ \ Ecuaciones exponenciales}}\\[3pt]
\mbox{\hyperlink{sec9}{\textbf{\color{fmgreen}2.9}\ \ Ecuaciones logarítmicas}}\\[3pt]\end{minipage}
\end{tcolorbox}

\Needspace{5.7cm}
\section{Polinomios. Operaciones}\hypertarget{sec1}{}
\begin{fmdef}{Definiciones}
\begin{itemize}
\item Un \textbf{monomio} es el producto de un número (\textbf{coeficiente}) y una o más letras elevadas a exponentes naturales (\textbf{parte literal}). El \textbf{grado} de un monomio es la suma de los exponentes de las letras. Por ejemplo, \fmm{-5x^3} tiene coeficiente \fmm{-5} y grado 3. Dos monomios son \textbf{semejantes} si tienen la misma parte literal; solo se pueden sumar o restar monomios semejantes.
\item Un \textbf{polinomio} en \fmm{x} es una suma de monomios: \fmm{\ P(x)=a_nx^n+a_{n-1}x^{n-1}+\cdots+a_1x+a_0,\ \ a_n\neq0.} Su \textbf{grado} es \fmm{n}; \fmm{a_n} es el \textbf{coeficiente principal} y \fmm{a_0}, el \textbf{término independiente}. Es \textbf{completo} si tiene todos los términos desde el grado \fmm{n} hasta el término independiente y está \textbf{ordenado} si los términos van de mayor a menor grado.
\item El \textbf{valor numérico} de \fmm{P(x)} para \fmm{x=a} es el número \fmm{P(a)} que se obtiene al sustituir \fmm{x} por \fmm{a}.
\item Dos polinomios son \textbf{iguales} si tienen el mismo grado y los coeficientes de los términos del mismo grado son iguales. Esto permite calcular coeficientes desconocidos \emph{igualando coeficientes}.
\end{itemize}
\end{fmdef}

\Needspace{5\baselineskip}
\begin{fmprop}{Operaciones con polinomios}
\begin{itemize}
\item \textbf{Suma y resta}: se suman (o restan) los monomios semejantes. \fmm{\ \text{grado}(P\pm Q)\le\max(\text{grado}\,P,\ \text{grado}\,Q)}.
\item \textbf{Producto}: se multiplica cada monomio del primero por cada monomio del segundo (propiedad distributiva) y se reducen los términos semejantes. \fmm{\ \text{grado}(P\cdot Q)=\text{grado}\,P+\text{grado}\,Q}.


\item \textbf{Identidades notables}

\begin{center}\renewcommand{\arraystretch}{0.4}\begin{tabularx}{\linewidth}{XX}
$\displaystyle (a+b)^2=a^2+2ab+b^2$ & $\displaystyle (a-b)^2=a^2-2ab+b^2$\\[3ex]
$\displaystyle (a+b)(a-b)=a^2-b^2$ & $\displaystyle (a+b+c)^2=a^2+b^2+c^2+2ab+2ac+2bc$\\[3ex]
$\displaystyle (a+b)^3=a^3+3a^2b+3ab^2+b^3$ & $\displaystyle (a-b)^3=a^3-3a^2b+3ab^2-b^3$\\[3ex]
\end{tabularx}\end{center}
\end{itemize}
\end{fmprop}

\Needspace{5\baselineskip}
\begin{fmex}{Ejemplo resuelto: operaciones con polinomios}
\begin{enumerate}[label=\textbf{\arabic*.},leftmargin=*]
\item Simplifica \fmm{3(x-1)^2-(2x+3)(2x-3)+x(x^2-4)}.

Desarrollamos con las identidades notables y la propiedad distributiva: \fmm{\ 3(x^2-2x+1)-(4x^2-9)+x^3-4x}

\fmm{=3x^2-6x+3-4x^2+9+x^3-4x=x^3-x^2-10x+12}
\item Calcula \fmm{a} y \fmm{b} para que \fmm{(x+a)(2x-3)=2x^2+bx-6}.

\fmm{(x+a)(2x-3)=2x^2+(2a-3)x-3a}

Igualamos coeficientes: \fmm{\ -3a=-6\Rightarrow a=2;\quad b=2a-3=1}
\end{enumerate}
\end{fmex}
\vspace*{0.3cm}
\Needspace{6.1cm}
\subsection*{División de polinomios}
\begin{fmdef}{División entera de polinomios}
Dados dos polinomios \fmm{D(x)} (\textbf{dividendo}) y \fmm{d(x)\neq0} (\textbf{divisor}), existen dos únicos polinomios \fmm{c(x)} (\textbf{cociente}) y \fmm{r(x)} (\textbf{resto}) tales que

\fmdisp{D(x)=d(x)\cdot c(x)+r(x)\qquad \text{con}\ \ \text{gr}\,r(x)\lt\text{gr}\,d(x)}

Si \fmm{r(x)=0}, la división es \textbf{exacta}: \fmm{d(x)} es un \textbf{divisor} de \fmm{D(x)}.

\begin{fmnote}\fmuline{Procedimiento}: se ordenan los polinomios (dejando hueco para los términos que faltan), se divide el primer término del dividendo entre el primero del divisor, se multiplica el resultado por el divisor y se resta. Se repite hasta que el grado de lo que queda es menor que el del divisor.\end{fmnote}
\end{fmdef}

\Needspace{5.7cm}
\section{Regla de Ruffini. Teorema del resto}\hypertarget{sec2}{}
\begin{fmdef}{Regla de Ruffini}
Cuando el divisor es de la forma \fmm{x-a}, la división se hace rápidamente con la \textbf{regla de Ruffini}:

\begin{enumerate}
\item Se escriben los coeficientes del dividendo, ordenados y \textbf{poniendo 0} en los términos que faltan. A la izquierda se escribe el número \fmm{a}.
\item Se baja el primer coeficiente. Se multiplica por \fmm{a}, se escribe debajo del coeficiente siguiente y se suman. Se repite hasta el final.
\item El último número es el \textbf{resto}; los anteriores son los coeficientes del \textbf{cociente}, que tiene un grado menos que el dividendo.
\end{enumerate}

\begin{fmnote}\fmuline{Atención}: si el divisor es \fmm{x+3}, entonces \fmm{a=-3}, porque \fmm{x+3=x-(-3)}.\end{fmnote}
\end{fmdef}

\Needspace{5\baselineskip}
\begin{fmprop}{Divisor de la forma $ax+b$}
Como \fmm{ax+b=a\left(x+\frac ba\right)}, hacemos Ruffini con el número \fmm{-\frac ba}. El \textbf{resto es el mismo}, pero el cociente obtenido se tiene que \textbf{dividir entre \fmm{a}}.

Por ejemplo, \fmm{(6x^3-x^2+4x-3):(2x-1)}: Ruffini con \fmm{\frac12} da \fmm{6x^2+2x+5} y resto \fmm{-\frac12}. Dividiendo entre 2: \fmm{\ c(x)=3x^2+x+\frac52,\ \ r=-\frac12}.
\end{fmprop}

\Needspace{5\baselineskip}
\begin{fmdef}{Teorema del resto y teorema del factor}
\textbf{Teorema del resto}: el resto de la división \fmm{P(x):(x-a)} es igual al valor numérico \fmm{P(a)}.

\fmdisp{r=P(a)}

\textbf{Demostración}: \fmm{P(x)=(x-a)\,c(x)+r}. Si sustituimos \fmm{x=a}: \fmm{P(a)=0\cdot c(a)+r=r}.

Un número \fmm{a} es una \textbf{raíz} (o cero) de \fmm{P(x)} si \fmm{P(a)=0}. \textbf{Teorema del factor}:

\fmdisp{P(a)=0\iff P(x)\ \text{es divisible por}\ (x-a)\iff P(x)=(x-a)\cdot c(x)}

\end{fmdef}

\Needspace{5\baselineskip}
\begin{fmex}{Ejemplo resuelto: teorema del resto}
\begin{enumerate}[label=\textbf{\arabic*.},leftmargin=*]
\item Calcula, sin hacer la división, el resto de \fmm{(x^5-3x^2+7):(x+1)}.

\fmm{r=P(-1)=(-1)^5-3(-1)^2+7=-1-3+7=3}
\item Halla \fmm{k} para que \fmm{P(x)=x^3+kx^2-5x+6} sea divisible por \fmm{x-3}.

Tiene que ser \fmm{P(3)=0}: \fmm{\ 27+9k-15+6=0}

\fmm{9k=-18\Rightarrow k=-2}
\end{enumerate}
\end{fmex}

\Needspace{5.7cm}
\section{Factorización de polinomios}\hypertarget{sec3}{}
\begin{fmdef}{Raíces y factorización}
\textbf{Factorizar} un polinomio es escribirlo como producto de polinomios del menor grado posible (\textbf{irreducibles}). Si \fmm{a_1,a_2,\ldots,a_n} son las raíces de \fmm{P(x)=a_nx^n+\cdots+a_0}, entonces

\fmdisp{P(x)=a_n\,(x-a_1)(x-a_2)\cdots(x-a_n)}

\begin{itemize}
\item Un polinomio de grado \fmm{n} tiene, como máximo, \fmm{n} raíces reales.
\item Si un factor \fmm{(x-a)} aparece \fmm{k} veces, \fmm{a} es una \textbf{raíz de multiplicidad \fmm{k}} (doble, triple\dots{}).
\item Las \textbf{raíces enteras} de un polinomio con coeficientes enteros son \textbf{divisores del término independiente} \fmm{a_0}. Las raíces racionales \fmm{\frac pq} (fracción irreducible) cumplen que \fmm{p} divide a \fmm{a_0} y \fmm{q} divide a \fmm{a_n}.
\item Hay polinomios sin raíces reales, como \fmm{x^2+1} o \fmm{x^2+x+1} \fmm{(\Delta\lt0)}: son irreducibles.
\end{itemize}
\end{fmdef}

\Needspace{5\baselineskip}
\begin{fmprop}{Procedimiento para factorizar}
\begin{enumerate}
\item \textbf{Sacar factor común} si es posible: \fmm{\ 2x^4-6x^3=2x^3(x-3)}.
\item Reconocer \textbf{identidades notables}: \fmm{\ x^2-9=(x+3)(x-3),\ \ 4x^2+4x+1=(2x+1)^2}.
\item Si es de \textbf{grado 2}, resolver la ecuación: \fmm{\ ax^2+bx+c=a(x-x_1)(x-x_2)}.
\item Si es \textbf{bicuadrado} \fmm{(ax^4+bx^2+c)}, hacer el cambio \fmm{t=x^2}.
\item Si es de grado mayor, aplicar \textbf{Ruffini} probando los divisores del término independiente (y, si hace falta, las fracciones \fmm{\frac pq}) hasta llegar a grado 2.
\end{enumerate}
\end{fmprop}

\Needspace{5\baselineskip}
\begin{fmex}{Ejemplo resuelto: factorización con Ruffini}
\begin{enumerate}[label=\textbf{\arabic*.},leftmargin=*]
\item Factoriza \fmm{P(x)=x^4-x^3-7x^2+x+6}.

Posibles raíces enteras: divisores de 6: \fmm{\ \pm1,\pm2,\pm3,\pm6}.

Ruffini con \fmm{x=1} (resto 0) y, con el cociente, \fmm{x=-1} (resto 0): \fmm{\ P(x)=(x-1)(x+1)(x^2-x-6)}

\fmm{x^2-x-6=0\Rightarrow x=\dfrac{1\pm5}{2}\Rightarrow x=3,\ x=-2}

\fmm{P(x)=(x-1)(x+1)(x-3)(x+2)}
\item Factoriza \fmm{Q(x)=2x^4+x^3-5x^2+2x}.

Factor común: \fmm{\ Q(x)=x(2x^3+x^2-5x+2)}

Ruffini con \fmm{x=1}: \fmm{\ 2x^3+x^2-5x+2=(x-1)(2x^2+3x-2)}

\fmm{2x^2+3x-2=0\Rightarrow x=\dfrac{-3\pm5}{4}\Rightarrow x=\frac12,\ x=-2}.\\
 Entonces \fmm{2x^2+3x-2=2\left(x-\frac12\right)(x+2)=(2x-1)(x+2) \Rightarrow Q(x)=x\,(x-1)(2x-1)(x+2)}
\end{enumerate}
\end{fmex}

\Needspace{5.7cm}
\section{Fracciones algebraicas}\hypertarget{sec4}{}
\begin{fmdef}{Definiciones}
\begin{itemize}
\item Una \textbf{fracción algebraica} es un cociente de dos polinomios, \fmm{\frac{P(x)}{Q(x)}}, con \fmm{Q(x)\neq0}. Solo tiene sentido para los valores de \fmm{x} que \textbf{no anulan el denominador}.
\item Dos fracciones \fmm{\frac PQ} y \fmm{\frac RS} son \textbf{equivalentes} si \fmm{P\cdot S=Q\cdot R}.
\item \textbf{Simplificar} es dividir numerador y denominador por un mismo factor: hay que \textbf{factorizar} y eliminar los factores comunes.
\item El \textbf{m.c.m.} de varios polinomios es el producto de los factores comunes y no comunes elevados al mayor exponente; el \textbf{m.c.d.}, el producto de los factores comunes elevados al menor exponente.
\end{itemize}

\begin{fmnote}\fmuline{Atención}: solo se pueden simplificar \emph{factores}, nunca \emph{sumandos}: \fmm{\ \dfrac{x+3}{x}\neq3}.\end{fmnote}
\end{fmdef}

\Needspace{5\baselineskip}
\begin{fmprop}{Operaciones con fracciones algebraicas}
\begin{itemize}
\item \textbf{Suma y resta}: se reducen a \textbf{común denominador} (m.c.m. de los denominadores) y se suman o restan los numeradores. \fmm{\ \dfrac PQ\pm\dfrac RS=\dfrac{P\cdot S\pm Q\cdot R}{Q\cdot S}}
\item \textbf{Producto}: \fmm{\ \dfrac PQ\cdot\dfrac RS=\dfrac{P\cdot R}{Q\cdot S}\qquad} \textbf{Cociente}: \fmm{\ \dfrac PQ:\dfrac RS=\dfrac{P\cdot S}{Q\cdot R}}
\item Es conveniente \textbf{factorizar antes de operar} y simplificar al final.
\end{itemize}
\end{fmprop}

\Needspace{5\baselineskip}
\begin{fmex}{Ejemplo resuelto: operaciones con fracciones}
\begin{enumerate}[label=\textbf{\arabic*.},leftmargin=*]
\item Simplifica \fmm{\dfrac{x^3-4x}{x^3-4x^2+4x}}.

\fmm{\dfrac{x(x^2-4)}{x(x^2-4x+4)}=\dfrac{x(x+2)(x-2)}{x(x-2)^2}=\dfrac{x+2}{x-2}}
\item Calcula \fmm{\dfrac{x}{x-1}-\dfrac{2}{x+1}+\dfrac{2}{x^2-1}}.

Factorizamos los denominadores: \fmm{\ x^2-1=(x+1)(x-1)\Rightarrow} m.c.m. \fmm{=(x+1)(x-1)}

\fmm{\dfrac{x(x+1)-2(x-1)+2}{(x+1)(x-1)}=\dfrac{x^2+x-2x+2+2}{(x+1)(x-1)}=\dfrac{x^2-x+4}{x^2-1}}
\newpage
\item Calcula y simplifica \fmm{\dfrac{x^2-9}{x^2+2x}\cdot\dfrac{x+2}{x-3}} y \fmm{\dfrac{2x-4}{x^2}:\dfrac{x^2-4}{x}}.

\fmm{\dfrac{(x+3)(x-3)}{x(x+2)}\cdot\dfrac{x+2}{x-3}=\dfrac{(x+3)(x-3)(x+2)}{x(x+2)(x-3)}=\dfrac{x+3}{x}}

\fmm{\dfrac{2(x-2)}{x^2}\cdot\dfrac{x}{(x+2)(x-2)}=\dfrac{2}{x(x+2)}}
\item Operaciones combinadas: calcula y simplifica \fmm{\left(\dfrac{x}{x-1}-\dfrac{1}{x}\right):\dfrac{x^2-x+1}{x^2-x}}.

Primero, el paréntesis (m.c.m. \fmm{=x(x-1)}): \fmm{\ \dfrac{x^2-(x-1)}{x(x-1)}=\dfrac{x^2-x+1}{x(x-1)}}

Después, el cociente (\fmm{x^2-x=x(x-1)}): \fmm{\ \dfrac{x^2-x+1}{x(x-1)}\cdot\dfrac{x(x-1)}{x^2-x+1}=1}
\end{enumerate}
\end{fmex}

\Needspace{8\baselineskip}
\section{Ecuaciones polinómicas}\hypertarget{sec5}{}

Una \textbf{ecuación polinómica} es una igualdad \fmm{P(x)=Q(x)} entre dos polinomios. El \textbf{grado} de la ecuación es el del polinomio que queda al pasarlo todo a un miembro y reducir. Las \textbf{soluciones} son los valores de \fmm{x} que hacen cierta la igualdad.

\Needspace{5\baselineskip}
\begin{fmdef}{Ecuaciones de primer grado}
Después de reducirla, una ecuación de primer grado tiene la forma \fmm{ax=b}. Para llegar a ella:

\begin{enumerate}
\item Si hay \textbf{denominadores}, se multiplican todos los términos (de los dos miembros) por el \textbf{m.c.m.} de los denominadores.
\item Se quitan los \textbf{paréntesis} aplicando la propiedad distributiva (cuidado con el signo \fmm{-} delante de un paréntesis o de una fracción: cambia el signo de \emph{todos} los términos).
\item Se \textbf{transponen} los términos: los que tienen \fmm{x} a un miembro y los números al otro, y se reducen.
\item Se \textbf{despeja} la \fmm{x}: \fmm{\ ax=b\Rightarrow x=\dfrac ba\ \ (a\neq0)}.
\end{enumerate}

\begin{fmnote}Si al reducir queda \fmm{0\cdot x=b} con \fmm{b\neq0}, la ecuación \textbf{no tiene solución}; si queda \fmm{0\cdot x=0}, es una \textbf{identidad} y cualquier número es solución.\end{fmnote}
\end{fmdef}

\Needspace{5\baselineskip}
\begin{fmex}{Ejemplo resuelto: ecuaciones de primer grado}
\begin{enumerate}[label=\textbf{\arabic*.},leftmargin=*]
\item Resuelve \fmm{3(2x-1)-2(x+4)=5-(x-2)}.

Quitamos los paréntesis: \fmm{\ 6x-3-2x-8=5-x+2}

Reducimos y transponemos: \fmm{\ 4x-11=7-x\Rightarrow4x+x=7+11\Rightarrow5x=18\Rightarrow x=\dfrac{18}{5}}
\item Resuelve \fmm{\dfrac{x+3}{4}-\dfrac{2x-1}{6}=\dfrac{x}{3}-1}.

m.c.m.\fmm{(4,6,3)=12}. Multiplicamos cada término por 12: \fmm{\ 3(x+3)-2(2x-1)=4x-12}

\fmm{3x+9-4x+2=4x-12\Rightarrow-x+11=4x-12\Rightarrow-5x=-23\Rightarrow x=\dfrac{23}{5}}

Fíjate en que el signo \fmm{-} delante de la segunda fracción cambia el signo de \fmm{2x} \emph{y} de \fmm{-1}.
\end{enumerate}
\end{fmex}

\Needspace{5\baselineskip}
\begin{fmdef}{Ecuaciones de segundo grado}
\begin{itemize}
\item \textbf{Segundo grado}: \fmm{\ ax^2+bx+c=0\ \Rightarrow\ x=\dfrac{-b\pm\sqrt{b^2-4ac}}{2a}}
\item El \textbf{discriminante} \fmm{\Delta=b^2-4ac} indica el número de soluciones: \fmm{\Delta\gt0}: dos soluciones; \fmm{\Delta=0}: una solución doble; \fmm{\Delta\lt0}: ninguna solución real.
\item \textbf{Incompletas}: \fmm{ax^2+bx=0\Rightarrow x(ax+b)=0\Rightarrow x=0,\ x=-\frac ba}; \fmm{\ ax^2+c=0\Rightarrow x=\pm\sqrt{-\frac ca}}.
\item Si \fmm{x_1} y \fmm{x_2} son las soluciones: \fmm{x_1+x_2=-\frac ba,\ \ x_1\cdot x_2=\frac ca}.
\end{itemize}
\end{fmdef}

\Needspace{5\baselineskip}
\begin{fmdef}{Ecuaciones bicuadradas y de grado superior}
\begin{itemize}
\item \textbf{Bicuadradas}: \fmm{ax^4+bx^2+c=0}. Se hace el cambio \fmm{t=x^2} y se obtiene \fmm{at^2+bt+c=0}. Por cada solución \fmm{t\ge0} se obtienen \fmm{x=\pm\sqrt t} (las \fmm{t\lt0} no dan soluciones reales).
\item El mismo método sirve para \fmm{ax^{2n}+bx^n+c=0} con \fmm{t=x^n} (por ejemplo, \fmm{x^6-9x^3+8=0} con \fmm{t=x^3}).
\item \textbf{Grado superior}: se pasa todo a un miembro, se \textbf{factoriza} (factor común, Ruffini\dots{}) y se aplica que un producto es cero si algún factor es cero:
\end{itemize}

\fmdisp{A(x)\cdot B(x)=0\iff A(x)=0\ \ \text{o}\ \ B(x)=0}

\end{fmdef}

\Needspace{5\baselineskip}
\begin{fmex}{Ejemplo resuelto: ecuaciones polinómicas}
\begin{enumerate}[label=\textbf{\arabic*.},leftmargin=*]
\item Resuelve \fmm{x^4-5x^2+4=0}.

\fmm{t=x^2:\ \ t^2-5t+4=0\Rightarrow t=\dfrac{5\pm3}{2}\Rightarrow t=4,\ t=1}

\fmm{x^2=4\Rightarrow x=\pm2;\quad x^2=1\Rightarrow x=\pm1.\quad x\in\{-2,-1,1,2\}}
\item Resuelve \fmm{x^3-2x^2-5x+6=0}.

Por Ruffini, \fmm{x=1} es raíz: \fmm{\ (x-1)(x^2-x-6)=0}

\fmm{(x-1)(x-3)(x+2)=0\Rightarrow x=1,\ x=3,\ x=-2}
\end{enumerate}
\end{fmex}

\Needspace{5.7cm}
\section{Ecuaciones con radicales}\hypertarget{sec6}{}
\begin{fmdef}{Ecuaciones irracionales}
Son ecuaciones en las que la incógnita aparece dentro de una raíz. Para resolverlas:

\begin{enumerate}
\item Se \textbf{aísla un radical} en uno de los miembros.
\item Se \textbf{elevan los dos miembros} al índice de la raíz (al cuadrado si es raíz cuadrada).
\item Si todavía queda algún radical, se repite el proceso.
\item Se resuelve la ecuación obtenida y \textbf{se comprueban todas las soluciones} en la ecuación inicial.
\end{enumerate}

\begin{fmnote}\fmuline{¿Por qué hay que comprobar?} Al elevar al cuadrado pueden aparecer \textbf{soluciones extrañas}, porque \fmm{A^2=B^2} también se cumple cuando \fmm{A=-B}. Con raíces de índice impar no aparecen, pero siempre es bueno comprobar.\end{fmnote}
\end{fmdef}

\Needspace{5\baselineskip}
\begin{fmex}{Ejemplo resuelto: ecuaciones con radicales}
\begin{enumerate}[label=\textbf{\arabic*.},leftmargin=*]
\item Resuelve \fmm{x-\sqrt{x+1}=5}.

Aislamos el radical: \fmm{\ x-5=\sqrt{x+1}}. Elevamos al cuadrado: \fmm{\ x^2-10x+25=x+1}

\fmm{x^2-11x+24=0\Rightarrow x=8,\ x=3}

Comprobación: \fmm{x=8:\ 8-\sqrt9=5\ \checkmark;\quad x=3:\ 3-\sqrt4=1\neq5\ \times}. Solución: \fmm{x=8}
\item Resuelve \fmm{\sqrt{2x+3}-\sqrt{x+1}=1}.

\fmm{\sqrt{2x+3}=1+\sqrt{x+1}\Rightarrow 2x+3=1+2\sqrt{x+1}+x+1}

\fmm{x+1=2\sqrt{x+1}\Rightarrow x^2+2x+1=4x+4\Rightarrow x^2-2x-3=0\Rightarrow x=3,\ x=-1}

Comprobación: \fmm{x=3:\ 3-2=1\ \checkmark;\quad x=-1:\ 1-0=1\ \checkmark}. Las dos son solución.
\end{enumerate}
\end{fmex}

\Needspace{5.7cm}
\section{Ecuaciones racionales}\hypertarget{sec7}{}
\begin{fmdef}{Ecuaciones racionales}
Son ecuaciones en las que la incógnita aparece en algún \textbf{denominador}. Para resolverlas:

\begin{enumerate}
\item Se factorizan los denominadores y se calcula el m.c.m.
\item Se multiplican los dos miembros por el m.c.m. para eliminar los denominadores.
\item Se resuelve la ecuación polinómica resultante.
\item \textbf{Se comprueban las soluciones}: hay que descartar las que anulan algún denominador.
\end{enumerate}
\end{fmdef}

\Needspace{5\baselineskip}
\begin{fmex}{Ejemplo resuelto: ecuaciones racionales}
\begin{enumerate}[label=\textbf{\arabic*.},leftmargin=*]
\item Resuelve \fmm{\dfrac{x}{x-1}+\dfrac{2}{x+1}=\dfrac{2}{x^2-1}}.

m.c.m. \fmm{=(x+1)(x-1)}. Multiplicamos los dos miembros: \fmm{\ x(x+1)+2(x-1)=2}

\fmm{x^2+3x-4=0\Rightarrow x=\dfrac{-3\pm5}{2}\Rightarrow x=1,\ x=-4}

Comprobación: \fmm{x=1} anula los denominadores \fmm{x-1} y \fmm{x^2-1}: se descarta. Solución: \fmm{x=-4}.
\item Resuelve \fmm{\dfrac{x}{x-2}-\dfrac{2}{x-2}=3}.

\fmm{x-2=3(x-2)\Rightarrow x-2=3x-6\Rightarrow x=2}

Pero \fmm{x=2} anula el denominador: la ecuación \textbf{no tiene solución}.
\end{enumerate}
\end{fmex}

\Needspace{5.7cm}
\section{Ecuaciones exponenciales}\hypertarget{sec8}{}
\begin{fmdef}{Ecuaciones exponenciales}
Son ecuaciones en las que la incógnita aparece en un \textbf{exponente}. Se basan en la propiedad:

\fmdisp{a^M=a^N\iff M=N\qquad(a\gt0,\ a\neq1)}

\textbf{Recuerda} las propiedades de las potencias: \fmm{\ a^m\cdot a^n=a^{m+n},\ \ \dfrac{a^m}{a^n}=a^{m-n},\ \ (a^m)^n=a^{m\cdot n},\ \ a^{-n}=\dfrac1{a^n},\ \ \sqrt[n]{a^m}=a^{m/n},\ \ a^0=1}.
\end{fmdef}

\Needspace{5\baselineskip}
\begin{fmprop}{Tipos de ecuaciones exponenciales y cómo resolverlas}
\begin{enumerate}
\item \textbf{Los dos miembros se pueden escribir como potencias de la misma base}: se igualan los exponentes.
\item \textbf{No se pueden reducir a la misma base}: se toman logaritmos en los dos miembros: \fmm{\ a^x=b\Rightarrow x=\dfrac{\log b}{\log a}}.
\item \textbf{Sumas de potencias de la misma base} con exponentes que difieren en un número: se descomponen (\fmm{2^{x+1}=2\cdot2^x,\ 2^{x-1}=\frac{2^x}{2}}) y se saca factor común \fmm{a^x}.
\item \textbf{Aparecen \fmm{a^{2x}} y \fmm{a^x}} (por ejemplo \fmm{4^x} y \fmm{2^x}): se hace el \textbf{cambio de variable} \fmm{t=a^x} (entonces \fmm{a^{2x}=t^2}) y se obtiene una ecuación de segundo grado. Como \fmm{t=a^x\gt0}, las soluciones \fmm{t\le0} no son válidas.
\end{enumerate}
\end{fmprop}

\Needspace{5\baselineskip}
\begin{fmex}{Ejemplo resuelto: uno de cada tipo}
\begin{enumerate}[label=\textbf{\arabic*.},leftmargin=*]
\item \fmm{8^{x+1}=4^{x^2}}

\fmm{2^{3x+3}=2^{2x^2}\Rightarrow 2x^2-3x-3=0\Rightarrow x=\dfrac{3\pm\sqrt{33}}{4}}
\newpage
\item \fmm{5^x=12}

\fmm{x\log5=\log12\Rightarrow x=\dfrac{\log12}{\log5}\approx 1{,}544}
\item \fmm{3^{x+1}+3^x+3^{x-1}=39}

\fmm{3^x\left(3+1+\tfrac13\right)=39\Rightarrow 3^x\cdot\dfrac{13}{3}=39}

\fmm{3^x=9\Rightarrow x=2}
\item \fmm{4^x-6\cdot2^x+8=0}

Como \fmm{4^x=(2^2)^x=(2^x)^2}, hacemos \fmm{t=2^x}: \fmm{\ t^2-6t+8=0\Rightarrow t=4,\ t=2}

\fmm{2^x=4\Rightarrow x=2;\quad 2^x=2\Rightarrow x=1}
\end{enumerate}
\end{fmex}

\Needspace{5.7cm}
\section{Ecuaciones logarítmicas}\hypertarget{sec9}{}
\begin{fmdef}{Logaritmos}
El \textbf{logaritmo} en base \fmm{a} \fmm{(a\gt0,\ a\neq1)} de un número \fmm{N\gt0} es el exponente al que hay que elevar \fmm{a} para obtener \fmm{N}:

\fmdisp{\log_aN=x\iff a^x=N}

\par\noindent \begin{minipage}[t]{0.48\linewidth}\vspace{0pt}\raggedright
\fmdisp{\begin{array}{l}\log_a(M\cdot N)=\log_aM+\log_aN\\[2pt] \log_a\dfrac MN=\log_aM-\log_aN\\[4pt] \log_aM^k=k\log_aM\end{array}}
\end{minipage}\hfill\begin{minipage}[t]{0.48\linewidth}\vspace{0pt}\raggedright
\fmdisp{\begin{array}{l}\log_aa=1\qquad\log_a1=0\\[4pt] \log_aN=\dfrac{\log N}{\log a}=\dfrac{\ln N}{\ln a}\end{array}}
\end{minipage}\par\medskip

\begin{fmnote}\fmuline{Atención}: solo existe el logaritmo de los números \textbf{positivos}.\end{fmnote}
\end{fmdef}

\Needspace{5\baselineskip}
\begin{fmprop}{Cómo resolver ecuaciones logarítmicas}
\begin{enumerate}
\item Usando las propiedades, se agrupan los logaritmos hasta llegar a \fmm{\log_aA=\log_aB}, y entonces \fmm{A=B}. Los números también se pueden escribir como logaritmos: \fmm{1=\log10,\ 2=\log100}\dots{}
\item O bien se llega a \fmm{\log_aA=k} y se aplica la definición: \fmm{A=a^k}.
\item \textbf{Se comprueban las soluciones}: todos los logaritmos de la ecuación original tienen que tener argumentos positivos.
\end{enumerate}
\end{fmprop}

\Needspace{5\baselineskip}
\begin{fmex}{Ejemplo resuelto: ecuaciones logarítmicas}
\begin{enumerate}[label=\textbf{\arabic*.},leftmargin=*]
\item Resuelve \fmm{\log x+\log(x-3)=1}.

\fmm{\log\big(x(x-3)\big)=\log10\Rightarrow x^2-3x-10=0\Rightarrow x=5,\ x=-2}

Comprobación: \fmm{x=-2} no es válida (\fmm{\log(-2)} no existe). Solución: \fmm{x=5}.
\item Resuelve \fmm{2\log x-\log(x+6)=0}.

\fmm{\log x^2=\log(x+6)\Rightarrow x^2-x-6=0\Rightarrow x=3,\ x=-2}

Solo es válida \fmm{x=3}.
\end{enumerate}
\end{fmex}

  
\end{document}
